\label{sec:experimental-design}

Empirical validation of the zone co-membership algorithm requires downloading
TIGER/Line and EIA Form 861 shapefiles, performing the spatial join to assign
zone membership to census tracts, and running \textsc{bisect} with the
\texttt{--weights-override zone-membership} flag.
These data are not yet downloaded and the implementation is not yet deployed;
this section documents the experimental design so that when the data becomes
available, the validation protocol is established.

\subsection{States Selected for Validation}

Three states are selected for initial validation based on the quality and coverage
of their zone data:

\begin{itemize}
  \item \textbf{North Carolina (NC)}: NC-14 (14 congressional districts) provides
    a well-studied redistricting case with extensive baseline results from earlier
    \textsc{bisect} runs.
    North Carolina has complete school district data in TIGER/Line (unified school
    districts covering all 100 counties) and good fire district coverage.
    NC is also a contested redistricting state, making it a high-priority validation
    target.
  \item \textbf{Wisconsin (WI)}: WI-8 (8 congressional districts) provides a
    mid-sized test case.
    Wisconsin has township-level county subdivision data (the town is a meaningful
    administrative unit in Wisconsin) and good fire district coverage.
  \item \textbf{Massachusetts (MA)}: MA-9 (9 congressional districts) provides the
    strongest test of the county subdivision dimension, because Massachusetts towns
    are the primary property tax-levying unit.
    Town integrity is a highly salient community criterion in Massachusetts, making
    it a natural ground truth for zone co-membership effects.
\end{itemize}

\subsection{Metrics}

\subsubsection{School District Preservation Rate}

The school district preservation rate is defined as the fraction of school district
pairs $(D_i, D_j)$ such that $D_i = D_j$ and both tracts end up in the same
redistricting district.
More precisely, for each pair of adjacent census tracts $(u, v)$ with
$\mathrm{zone}_{\text{school}}(u) = \mathrm{zone}_{\text{school}}(v)$, the pair
is \emph{preserved} if $\mathrm{district}(u) = \mathrm{district}(v)$.
The preservation rate is the fraction of such adjacent co-school-district pairs that
are preserved.

A preservation rate of 100\% would mean that no two tracts in the same school district
were split across different redistricting districts.
A preservation rate of 0\% would mean that every adjacent co-school-district pair
was split.
Geographic baseline (standard-bisect) preservation rates for NC, WI, and MA will be
measured first; the zone co-membership mode should produce strictly higher preservation
rates.

\subsubsection{School District Boundary Crossing Count}

The school district boundary crossing count is the total number of edges $(u, v) \in E$
in the final partition where:
\begin{enumerate}
  \item Tracts $u$ and $v$ are assigned to different redistricting districts, AND
  \item $\mathrm{zone}_{\text{school}}(u) \neq \mathrm{zone}_{\text{school}}(v)$.
\end{enumerate}
This metric counts cuts at school district boundaries specifically.
Under the L0 test invariant \texttt{school\_district\_split\_never\_increases},
this count must be $\leq$ the count produced by geographic weights.
The empirical claim in the M.6 spec is that the count is $\leq 20\%$ of the geographic
weight count for NC-14.

\subsubsection{Compactness Change}

Zone co-membership weights may trade compactness for administrative coherence.
If zone boundaries follow roads and natural features (as they typically do), zone
co-membership cuts should be at least as compact as geographic cuts.
If zone boundaries cross irregular terrain, zone co-membership weights may produce
slightly less compact districts.
The Polsby-Popper and Reock compactness scores will be computed for all districts
under both weight modes (geographic baseline and zone-membership), and the mean
and standard deviation of the change will be reported.
The hypothesis is that zone-membership weights produce no statistically significant
compactness degradation, because administrative boundaries tend to follow the same
roads and rivers that are natural shape boundaries.

\subsubsection{Proportionality Gap}

The proportionality gap measures the difference between a party's expected seat share
and its statewide vote share.
Zone co-membership weights are partisan-neutral by construction, but they may
nonetheless affect the proportionality gap if partisan registration is geographically
correlated with administrative zone membership (e.g., if school district boundaries
align with partisan patterns).
The proportionality gap will be computed for NC, WI, and MA under both weight modes
to document whether zone co-membership weights introduce or reduce partisan bias
compared to the geographic baseline.
No hypothesis is stated for this metric: it is reported for transparency, not as a
validation criterion.

\subsection{Comparison Protocol}

All experiments run using the same configuration except for the weight mode:
\begin{verbatim}
bisect state --state NC --year 2020 \
  --structure prime-factor \
  --search convergence \
  --weights-override geographic \
  --version v_geo_baseline

bisect state --state NC --year 2020 \
  --structure prime-factor \
  --search convergence \
  --weights-override zone-membership \
  --version v_zone_m6
\end{verbatim}

The convergence search strategy (seed search until $T = 600$ seconds) is used
for both to ensure that the comparison is not confounded by seed selection.
Both runs use the 2020 census year and the prime-factor bisection structure,
consistent with the platform defaults.

\subsection{New England Town Integrity}

For Massachusetts, an additional metric specific to New England is computed:
\textbf{town integrity}, defined as the fraction of Massachusetts towns that are
entirely contained within a single redistricting district.
Under geographic weights, small towns in densely populated areas are frequently
split across districts.
Under zone co-membership weights, the county subdivision dimension should strongly
incentivize keeping towns whole.
The target from the M.6 spec is that $\geq 95\%$ of Massachusetts towns remain
whole in $\geq 95\%$ of zone-membership runs.

This metric has direct legal relevance for Massachusetts redistricting: the
Massachusetts Constitution (Part II, Ch. 1, \S3) has historically been interpreted
to require preservation of town boundaries in legislative redistricting.
The town integrity metric thus doubles as a statutory compliance check for
Massachusetts legislative district plans.
